\documentclass{zhvt-classic}

% 使用本地源流明体（GenRyuMin2 TW）按文件路径加载，不依赖系统字体名
% 版框（内框／界栏／书口侧栏／鱼尾）改用与版心印章同一印泥棕色：
%   印章图片 yshs.png 实色 = RGB(135,68,52) = HTML 874434
\usepackage{xcolor}
\usepackage{tikz}
% sealbrown 双注册：xcolor 供 TikZ/pgf，L3 供 cls 版框
\definecolor{sealbrown}{HTML}{874434}
\ExplSyntaxOn
\color_set:nnn { sealbrown } { HTML } { 874434 }
\ExplSyntaxOff

\zhvtset{
  grid color       = sealbrown,
  main font        = GenRyuMin2TW-SB,
  main font path   = C:/Users/Administrator/texmf/fonts/opentype/public/GenRyuMin2TW/,
jiazhu font      = GenRyuMin2TW-M,
  jiazhu font path = C:/Users/Administrator/texmf/fonts/opentype/public/GenRyuMin2TW/,
  % 页码框容纳 4 字：正常页「九十七」+前置空 1 格=4 行恰好；
  % 100+ 页「一百零一」共 5 行超出，末字「一/二/三」被框裁掉（2026-10-04 实测），故放宽到 5
  shukou page number chars = 5,
  % 书口书名框高从默认 6 字减到 4 字，把省出的 2 字高拨给 chapter 框，
  % 使 6 字书名+卷名放进 9 字高的 chapter 框后不再压线（2026-10-05）
  shukou title chars = 4,
}

\title{史记}[酷吏列传]
\author{汉·司马迁}
\maker{闲叟漫记·重製}


% ===== 索隐述赞样式（2026-11-08 四次定稿 · 愚大人亲手改定参数）=====
% ⚠ 竖排前提：自定义字体家族必须带 RawFeature={vertical:+vert}；
%   simhei.ttf 加 +vhal 会整体加载失败静默回退，只写 +vert；simfang 同理。
\newCJKfontfamily{\shuzanfont}[Path=C:/Windows/Fonts/,RawFeature={vertical:+vert}]{simhei.ttf}
\newCJKfontfamily{\shuzanbody}[Path=C:/Windows/Fonts/,RawFeature={vertical:+vert}]{simfang.ttf}
% 题头：152x30pt 圆角 6pt 棕底白字；baseline (0,2) 使整框落入界行格子不压版框线；
% 白字节点 x 不补偿、y +8.2pt（像素级墨迹实测 dx/dy≈0，L/R 留白约 20pt）。
\newcommand{\shuzanhead}[1]{%
  \begin{tikzpicture}[baseline={(0,2)}]
    \fill[sealbrown, rounded corners=6pt]
      (-76pt,-15pt) rectangle (76pt,15pt);
    \node[text=white, inner sep=0pt, anchor=center] at (0pt,8.2pt)
      {\shuzanfont\fontsize{23}{23}\selectfont #1};
  \end{tikzpicture}%
}
% ★ 修正书口卷名位置：chapter 框起始位置改回 cls 默认值，避免全局影响。
%   只在每个卷首页通过 \juan 命令里的空命令微调该页位置（安全，不会跨页污染）。
\ExplSyntaxOn
\AtBeginDocument{
  \dim_set:Nn \l__zhvt_shukou_chapter_title_start_dim
    { \l__zhvt_shukou_title_height_dim + \l__zhvt_font_size_dim * 2 }
  \dim_set:Nn \l__zhvt_shukou_page_number_start_dim
    { \l__zhvt_shukou_chapter_title_start_dim + \l__zhvt_shukou_chapter_title_height_dim }
}
\ExplSyntaxOff

\begin{document}

% ★★ 版心（书口）排布覆写：把书名放到鱼尾之下（2026-10-04）
%
%   cls 默认版心自上而下：书名（title 静态框，在鱼尾之上）→ 鱼尾 → 篇名（chapter 动态框）→ 页码
%   本模板改为：          鱼尾（位置不动）→ 书名 → 篇名 → 页码分割线 → 页码
%
%   手法：把 cls 放在 title 静态框里的书名，改放进 chapter 动态框（鱼尾之下）；
%         title 静态框留空。书名取 \l__zhvt_title_str，即 \title{...} 的第一参数，
%         所以换书时只改 \title，不必在此重复写书名。
%
%   ⚠ 三条必须遵守的写法（2026-10-04 实测）：
%     ① 汉字必须用 \hbox{...} 包住。flowfram 会把框内容装进一个宽 14.4mm（= 一个 line_sep
%        = 书口栏宽）的纵向 minipage 并按该宽度折行，而整页靠 PDF /Rotate 90 显示、
%        框内文字本身不变向 —— 不包 \hbox 时「围炉夜话」4 字会被逐字折成 4 行，
%        在版面上就成了横排 4 字并横向溢出版框。
%     ② 竖排下**一个 \hbox 就是一整列**。书名与篇名要同处一列，必须放进同一个 \hbox；
%        分成两个 \hbox 时，后一个会被当成新的一列、甩到隔壁栏去
%        （实测「序」跑到书口栏外侧 14.4mm 的正文列上）。中间的 \quad 在框内是横向空白，
%        显示出来正好是书名与篇名之间的纵向间隔。
%     ③ ★ 内容盒要写成 **与框等宽的定宽盒** `\hbox to \l__zhvt_shukou_chapter_title_height_dim{…\hfil}`。
%        实测：chapter 框是把内容盒沿纵向**居中**摆放的，内容宽度一变，书名就跟着漂：
%        自然宽度的 \hbox（书名 38.4mm）落在 x=157.2mm，加篇名后（57.6mm）落在 x=152.4mm，
%        目录页（无篇名）与正文页的书名就差半格。做成等宽盒后，无论有没有篇名，
%        内容固定从框顶起排，实测两页均落在 x=142.8mm，不再漂。
\ExplSyntaxOn
% 把书口 chapter 框（书名+篇名）起始位置下移 1 字高，避免 4 字卷名顶到鱼尾
% ⚠ cls 的 \__zhvt_calculate_dimensions: 在 \AtBeginDocument 里执行，会重算该 dim；
%   必须把覆写挂到 begin document 之后，否则会被二次覆盖（2026-10-05 实测）。
% (旧方案已移到导言区，直接修改 chapter 框起始位置)
% 卷标题命令 #1=书口篇名(鱼尾) #2=正文卷名（含卷序，卷名内空格用 ~）
%
% ★★ 目录页两项修订（2026-10-05，愚大人指出）★★
%   ① 页码起点统一对齐：
%      原状 —— 目录项写成 \par\hspace{4em}卷名\hspace{4em}页码，页码起点 = 卷名实宽 + 4em。
%      卷名实宽随字数变：「卷一~梁惠王上」6 字 +~、「卷七~离娄上」5 字 +~，
%      差 1 个字位 → 卷七~卷十 4 条的页码整体比其余 10 条高 27.21pt（一个字），看着上下不齐。
%      改法 —— 卷名补全角空格（宽 = 1 字 = 27.215pt）补足到 **7 字符位**再写进目录，
%      14 条页码起点遂全部落在同一条竖线上（起点对齐；末字随页码字数自然参差）。
%      ⚠ 只对目录里的副本 \l_tmpc_tl 补，正文卷名仍用 \l_tmpa_tl 原样 —— 正文像素不受影响。
%   ② 卷名整体上移半字位：前导 \hspace 由 4em 改 3.5em。
%      实测原状目录块在版框内偏下（天头留白 139.6pt / 地脚留白 110.7pt），
%      上移 0.5em = 13.6pt 后上下基本均分，视觉居中。
\NewDocumentCommand { \juan } { mm } {
	\clearpage
	\tl_set:Nn \l__zhvt_chapter_title_tl { #1 }
	\tl_set:Nn \l_tmpa_tl { #2 }
	\tl_set_eq:NN \l_tmpc_tl \l_tmpa_tl
	\int_step_inline:nn { 7 - \str_count:n { #2 } }
	  { \tl_put_right:Nn \l_tmpc_tl { 　 } }
	\tl_set:Nx \l_tmpb_tl { \zhnum { page } }
	\tl_put_left:Nn \l_tmpb_tl { \hspace { 4em } }
	\tl_put_left:Nn \l_tmpb_tl { \l_tmpc_tl }
	\tl_put_left:Nn \l_tmpb_tl { \par \hspace { 3.5em } }
	\addtocontents { toc } { \l_tmpb_tl }
	\skip_horizontal:n { \l__zhvt_font_size_dim * 4 }
	\l_tmpa_tl
	\par
	\thispagestyle { empty }
}
\cs_set_protected:Npn \__zhvt_set_shukou_content: {
	% ★ 鱼尾上方留白（title 静态框）嵌入闲叟钤印：
	%   整页 /Rotate 90，故此处把图用 \rotatebox{90} 转正为竖章；
	%   尺寸卡在版心条幅内（书内约 7mm 宽 × 33.8mm 高，留余量不压鱼尾）。
	\setstaticcontents * { title } {
		\vbox_to_ht:nn { \l__zhvt_baseline_skip_dim } {
			\vskip 14.2pt
			\hbox_to_wd:nn { \l__zhvt_shukou_title_height_dim } {
				\hfil
				\rotatebox [ origin = c ] { 90 } { \includegraphics [ width = 7mm , height = 33.8mm ] { yshs.png } }
				\hfil
			}
		}
	}
% ★ 版心（书口）统一改法（2026-11-16 愚大人定）：18pt＋框内居中＋书口文字同版框色
%   ① 原 27.2pt 时 7 字长篇名（如「廉颇蔺相如列传」）末字溢出 chapter 框、压住页码；
%      改 \fontsize{18pt}{18pt}，前后 \hfil 令字串在框内居中（不随篇名字数漂移）。
%   ② \vbox{\vskip 8.2pt} 定量下移，抵消缩字号后栏宽方向 8.2pt 上偏（像素实测 y 中心 543.3 = 原版）。
%      ⚠ \raisebox / \vspace* 在 flowfram 框内无效；换字号须同步换补偿：16pt→10.1、19pt→7.4。
%   ③ chapter 与 page~number 均置 \color{sealbrown}（HTML 874434），与版框／界行同色；正文仍黑。
%   ⚠ 页码字位、钤印、鱼尾、界行、正文一律未动；五部合订本同此口径。
	\setdynamiccontents * { chapter }
	{ \color{sealbrown} \vbox { \vskip 8.2pt \hbox to \l__zhvt_shukou_chapter_title_height_dim
		{ \hfil \fontsize{18pt}{18pt}\selectfont \l__zhvt_title_str \quad \l__zhvt_chapter_title_tl \hfil } } }
	\setdynamiccontents * { page~number } { \color{sealbrown} \quad \zhnum { page } }
}
\__zhvt_set_shukou_content:
\ExplSyntaxOff

\maketitle

% 单篇，无多卷目录，不渲染目录页

\juan{酷吏列传}{卷一百二十二~酷吏列传第六十二}
\hspace*{2em}孔子曰。﹁道之以政。齐之以刑。民免而无耻。道之以德。齐之以礼。有耻且格。﹂老氏称。﹁上德不德。是以有德。下德不失德。是以无德。法令滋章。盗贼多有。﹂太史公曰。信哉是言也。法令者治之具。而非制治清浊之源也。昔天下之网尝密矣。然奸伪萌起。其极也。上下相遁。至於不振。当是之时。吏治若救火扬沸。非武健严酷。恶能胜其任而愉快乎。言道德者。溺其职矣。故曰﹁听讼。吾犹人也。必也使无讼乎﹂。﹁下士闻道大笑之﹂。非虚言也。汉兴。破觚而为圜。斫雕而为朴。网漏於吞舟之鱼。而吏治烝烝。不至於奸。黎民艾安。由是观之。在彼不在此。\\
\hspace*{2em}高后时。酷吏独有侯封。刻轹宗室。侵辱功臣。吕氏已败。遂夷侯封之家。孝景时。晁错以刻深颇用术辅其资。而七国之乱。发怒於错。错卒以被戮。其後有郅都。宁成之属。\\
\hspace*{2em}郅都者。杨人也。以郎事孝文帝。孝景时。都为中郎将。敢直谏。面折大臣於朝。尝从入上林。贾姬如厕。野彘卒入厕。上目都。都不行。上欲自持兵救贾姬。都伏上前曰。﹁亡一姬复一姬进。天下所少宁贾姬等乎。陛下纵自轻。柰宗庙太后何。﹂上还。彘亦去。太后闻之。赐都金百斤。由此重郅都。\\
\hspace*{2em}济南瞷氏宗人三百馀家。豪猾。二千石莫能制。於是景帝乃拜都为济南太守。至则族灭瞷氏首恶。馀皆股栗。居岁馀。郡中不拾遗。旁十馀郡守畏都如大府。\\
\hspace*{2em}都为人勇。有气力。公廉。不发私书。问遗无所受。请寄无所听。常自称曰。﹁已倍亲而仕。身固当奉职死节官下。终不顾妻子矣。﹂\\
\hspace*{2em}郅都迁为中尉。丞相条侯至贵倨也。而都揖丞相。是时民朴。畏罪自重。而都独先严酷。致行法不避贵戚。列侯宗室见都侧目而视。号曰﹁苍鹰﹂。\\
\hspace*{2em}临江王徵诣中尉府对簿。临江王欲得刀笔为书谢上。而都禁吏不予。魏其侯使人以间与临江王。临江王既为书谢上。因自杀。窦太后闻之。怒。以危法中都。都免归家。孝景帝乃使使持节拜都为雁门太守。而便道之官。得以便宜从事。匈奴素闻郅都节。居边。为引兵去。竟郅都死不近雁门。匈奴至为偶人象郅都。令骑驰射莫能中。见惮如此。匈奴患之。窦太后乃竟中都以汉法。景帝曰。﹁都忠臣。﹂欲释之。窦太后曰。﹁临江王独非忠臣邪。﹂於是遂斩郅都。\\
\hspace*{2em}宁成者。穰人也。以郎谒者事景帝。好气。为人小吏。必陵其长吏。为人上。操下如束湿薪。滑贼任威。稍迁至济南都尉。而郅都为守。始前数都尉皆步入府。因吏谒守如县令。其畏郅都如此。及成往。直陵都出其上。都素闻其声。於是善遇。与结驩。久之。郅都死。後长安左右宗室多暴犯法。於是上召宁成为中尉。其治效郅都。其廉弗如。然宗室豪桀皆人人惴恐。\\
\hspace*{2em}武帝即位。徙为内史。外戚多毁成之短。抵罪髡钳。是时九卿罪死即死。少被刑。而成极刑。自以为不复收。於是解脱。诈刻传出关归家。称曰。﹁仕不至二千石。贾不至千万。安可比人乎。﹂乃贳贷买陂田千馀顷。假贫民。役使数千家。数年。会赦。致产数千金。为任侠。持吏长短。出从数十骑。其使民威重於郡守。\\
\hspace*{2em}周阳由者。其父赵兼以淮南王舅父侯周阳。故因姓周阳氏。由以宗家任为郎。事孝文及景帝。景帝时。由为郡守。武帝即位。吏治尚循谨甚。然由居二千石中。最为暴酷骄恣。所爱者。挠法活之。所憎者。曲法诛灭之。所居郡。必夷其豪。为守。视都尉如令。为都尉。必陵太守。夺之治。与汲黯俱为忮。司马安之文恶。俱在二千石列。同车未尝敢均茵伏。\\
\hspace*{2em}由後为河东都尉。时与其守胜屠公争权。相告言罪。胜屠公当抵罪。义不受刑。自杀。而由弃市。\\
\hspace*{2em}自宁成。周阳由之後。事益多。民巧法。大抵吏之治类多成。由等矣。\\
\hspace*{2em}赵禹者。斄人。以佐史补中都官。用廉为令史。事太尉亚夫。亚夫为丞相。禹为丞相史。府中皆称其廉平。然亚夫弗任。曰。﹁极知禹无害。然文深。不可以居大府。﹂今上时。禹以刀笔吏积劳。稍迁为御史。上以为能。至太中大夫。与张汤论定诸律令。作见知。吏传得相监司。用法益刻。盖自此始。\\
\hspace*{2em}张汤者。杜人也。其父为长安丞。出。汤为兒守舍。还而鼠盗肉。其父怒。笞汤。汤掘窟得盗鼠及馀肉。劾鼠掠治。传爰书。讯鞫论报。并取鼠与肉。具狱磔堂下。其父见之。视其文辞如老狱吏。大惊。遂使书狱。父死後。汤为长安吏。久之。\\
\hspace*{2em}周阳侯始为诸卿时。尝系长安。汤倾身为之。及出为侯。大与汤交。遍见汤贵人。汤给事内史。为宁成掾。以汤为无害。言大府。调为茂陵尉。治方中。\\
\hspace*{2em}武安侯为丞相。徵汤为史。时荐言之天子。补御史。使案事。治陈皇后蛊狱。深竟党与。於是上以为能。稍迁至太中大夫。与赵禹共定诸律令。务在深文。拘守职之吏。已而赵禹迁为中尉。徙为少府。而张汤为廷尉。两人交驩。而兄事禹。禹为人廉倨。为吏以来。舍毋食客。公卿相造请禹。禹终不报谢。务在绝知友宾客之请。孤立行一意而已。见文法辄取。亦不覆案。求官属阴罪。汤为人多诈。舞智以御人。始为小吏。乾没。与长安富贾田甲。鱼翁叔之属交私。及列九卿。收接天下名士大夫。己心内虽不合。然阳浮慕之。\\
\hspace*{2em}是时上方乡文学。汤决大狱。欲傅古义。乃请博士弟子治尚书。春秋补廷尉史。亭疑法。奏谳疑事。必豫先为上分别其原。上所是。受而著谳决法廷尉。絜令扬主之明。奏事即谴。汤应谢。乡上意所便。必引正。监。掾史贤者。曰。﹁固为臣议。如上责臣。臣弗用。愚抵於此。﹂罪常释。即奏事。上善之。曰。﹁臣非知为此奏。乃正。监。掾史某为之。﹂其欲荐吏。扬人之善蔽人之过如此。所治即上意所欲罪。予监史深祸者。即上意所欲释。与监史轻平者。所治即豪。必舞文巧诋。即下户羸弱。时口言。虽文致法。上财察。於是往往释汤所言。汤至於大吏。内行脩也。通宾客饮食。於故人子弟为吏及贫昆弟。调护之尤厚。其造请诸公。不避寒暑。是以汤虽文深意忌不专平。然得此声誉。而刻深吏多为爪牙用者。依於文学之士。丞相弘数称其美。及治淮南。衡山。江都反狱。皆穷根本。严助及伍被。上欲释之。汤争曰。﹁伍被本画反谋。而助亲幸出入禁闼爪牙臣。乃交私诸侯如此。弗诛。後不可治。﹂於是上可论之。其治狱所排大臣自为功。多此类。於是汤益尊任。迁为御史大夫。\\
\hspace*{2em}会浑邪等降。汉大兴兵伐匈奴。山东水旱。贫民流徙。皆仰给县官。县官空虚。於是丞上指。请造白金及五铢钱。笼天下盐铁。排富商大贾。出告缗令。鉏豪彊并兼之家。舞文巧诋以辅法。汤每朝奏事。语国家用。日晏。天子忘食。丞相取充位。天下事皆决於汤。百姓不安其生。骚动。县官所兴。未获其利。奸吏并侵渔。於是痛绳以罪。则自公卿以下。至於庶人。咸指汤。汤尝病。天子至自视病。其隆贵如此。\\
\hspace*{2em}匈奴来请和亲。群臣议上前。博士狄山曰。﹁和亲便。﹂上问其便。山曰。﹁兵者凶器。未易数动。高帝欲伐匈奴。大困平城。乃遂结和亲。孝惠。高后时。天下安乐。及孝文帝欲事匈奴。北边萧然苦兵矣。孝景时。吴楚七国反。景帝往来两宫间。寒心者数月。吴楚已破。竟景帝不言兵。天下富实。今自陛下举兵击匈奴。中国以空虚。边民大困贫。由此观之。不如和亲。﹂上问汤。汤曰。﹁此愚儒。无知。﹂狄山曰。﹁臣固愚忠。若御史大夫汤乃诈忠。若汤之治淮南。江都。以深文痛诋诸侯。别疏骨肉。使蕃臣不自安。臣固知汤之为诈忠。﹂於是上作色曰。﹁吾使生居一郡。能无使虏入盗乎。﹂曰。﹁不能。﹂曰。﹁居一县。﹂对曰。﹁不能。﹂复曰。﹁居一障间。﹂山自度辩穷且下吏。曰。﹁能。﹂於是上遣山乘鄣。至月馀。匈奴斩山头而去。自是以後。群臣震慴。\\
\hspace*{2em}汤之客田甲。虽贾人。有贤操。始汤为小吏时。与钱通。及汤为大吏。甲所以责汤行义过失。亦有烈士风。\\
\hspace*{2em}汤为御史大夫七岁。败。\\
\hspace*{2em}河东人李文尝与汤有卻。已而为御史中丞。恚。数从中文书事有可以伤汤者。不能为地。汤有所爱史鲁谒居。知汤不平。使人上蜚变告文奸事。事下汤。汤治论杀文。而汤心知谒居为之。上问曰。﹁言变事纵迹安起。﹂汤详惊曰。﹁此殆文故人怨之。﹂谒居病卧闾里主人。汤自往视疾。为谒居摩足。赵国以冶铸为业。王数讼铁官事。汤常排赵王。赵王求汤阴事。谒居尝案赵王。赵王怨之。并上书告。﹁汤。大臣也。史谒居有病。汤至为摩足。疑与为大奸。﹂事下廷尉。谒居病死。事连其弟。弟系导官。汤亦治他囚导官。见谒居弟。欲阴为之。而详不省。谒居弟弗知。怨汤。使人上书告汤与谒居谋。共变告李文。事下减宣。宣尝与汤有卻。及得此事。穷竟其事。未奏也。会人有盗发孝文园瘗钱。丞相青翟朝。与汤约俱谢。至前。汤念独丞相以四时行园。当谢。汤无与也。不谢。丞相谢。上使御史案其事。汤欲致其文丞相见知。丞相患之。三长史皆害汤。欲陷之。\\
\hspace*{2em}始长史硃买臣。会稽人也。读春秋。庄助使人言买臣。买臣以楚辞与助俱幸。侍中。为太中大夫。用事。而汤乃为小吏。跪伏使买臣等前。已而汤为廷尉。治淮南狱。排挤庄助。买臣固心望。及汤为御史大夫。买臣以会稽守为主爵都尉。列於九卿。数年。坐法废。守长史。见汤。汤坐床上。丞史遇买臣弗为礼。买臣楚士。深怨。常欲死之。王朝。齐人也。以术至右内史。边通。学长短。刚暴彊人也。官再至济南相。故皆居汤右。已而失官。守长史。诎体於汤。汤数行丞相事。知此三长史素贵。常凌折之。以故三长史合谋曰。﹁始汤约与君谢。已而卖君。今欲劾君以宗庙事。此欲代君耳。吾知汤阴事。﹂使吏捕案汤左田信等。曰汤且欲奏请。信辄先知之。居物致富。与汤分之。及他奸事。事辞颇闻。上问汤曰。﹁吾所为。贾人辄先知之。益居其物。是类有以吾谋告之者。﹂汤不谢。汤又详惊曰。﹁固宜有。﹂减宣亦奏谒居等事。天子果以汤怀诈面欺。使使八辈簿责汤。汤具自道无此。不服。於是上使赵禹责汤。禹至。让汤曰。﹁君何不知分也。君所治夷灭者几何人矣。今人言君皆有状。天子重致君狱。欲令君自为计。何多以对簿为。﹂汤乃为书谢曰。﹁汤无尺寸功。起刀笔吏。陛下幸致为三公。无以塞责。然谋陷汤罪者。三长史也。﹂遂自杀。\\
\hspace*{2em}汤死。家产直不过五百金。皆所得奉赐。无他业。昆弟诸子欲厚葬汤。汤母曰。﹁汤为天子大臣。被汙恶言而死。何厚葬乎。﹂载以牛车。有棺无椁。天子闻之。曰。﹁非此母不能生此子。﹂乃尽案诛三长史。丞相青翟自杀。出田信。上惜汤。稍迁其子安世。\\
\hspace*{2em}赵禹中废。已而为廷尉。始条侯以为禹贼深。弗任。及禹为少府。比九卿。禹酷急。至晚节。事益多。吏务为严峻。而禹治加缓。而名为平。王温舒等後起。治酷於禹。禹以老。徙为燕相。数岁。乱悖有罪。免归。後汤十馀年。以寿卒于家。\\
\hspace*{2em}义纵者。河东人也。为少年时。尝与张次公俱攻剽为群盗。纵有姊姁。以医幸王太后。王太后问。﹁有子兄弟为官者乎。﹂姊曰。﹁有弟无行。不可。﹂太后乃告上。拜义姁弟纵为中郎。补上党郡中令。治敢行。少蕴藉。县无逋事。举为第一。迁为长陵及长安令。直法行治。不避贵戚。以捕案太后外孙脩成君子仲。上以为能。迁为河内都尉。至则族灭其豪穰氏之属。河内道不拾遗。而张次公亦为郎。以勇悍从军。敢深入。有功。为岸头侯。\\
\hspace*{2em}宁成家居。上欲以为郡守。御史大夫弘曰。﹁臣居山东为小吏时。宁成为济南都尉。其治如狼牧羊。成不可使治民。﹂上乃拜成为关都尉。岁馀。关东吏隶郡国出入关者。号曰﹁宁见乳虎。无值宁成之怒﹂。义纵自河内迁为南阳太守。闻宁成家居南阳。及纵至关。宁成侧行送迎。然纵气盛。弗为礼。至郡。遂案宁氏。尽破碎其家。成坐有罪。及孔。暴之属皆饹亡。南阳吏民重足一迹。而平氏硃彊。杜衍。杜周为纵牙爪之吏。任用。迁为廷史。军数出定襄。定襄吏民乱败。於是徙纵为定襄太守。纵至。掩定襄狱中重罪轻系二百馀人。及宾客昆弟私入相视亦二百馀人。纵一捕鞠。曰﹁为死罪解脱﹂。是日皆报杀四百馀人。其後郡中不寒而栗。猾民佐吏为治。\\
\hspace*{2em}是时赵禹。张汤以深刻为九卿矣。然其治尚宽。辅法而行。而纵以鹰击毛挚为治。後会五铢钱白金起。民为奸。京师尤甚。乃以纵为右内史。王温舒为中尉。温舒至恶。其所为不先言纵。纵必以气凌之。败坏其功。其治。所诛杀甚多。然取为小治。奸益不胜。直指始出矣。吏之治以斩杀缚束为务。阎奉以恶用矣。纵廉。其治放郅都。上幸鼎湖。病久。已而卒起幸甘泉。道多不治。上怒曰。﹁纵以我为不复行此道乎。﹂嗛之。至冬。杨可方受告缗。纵以为此乱民。部吏捕其为可使者。天子闻。使杜式治。以为废格沮事。弃纵市。後一岁。张汤亦死。\\
\hspace*{2em}王温舒者。阳陵人也。少时椎埋为奸。已而试补县亭长。数废。为吏。以治狱至廷史。事张汤。迁为御史。督盗贼。杀伤甚多。稍迁至广平都尉。择郡中豪敢任吏十馀人。以为爪牙。皆把其阴重罪。而纵使督盗贼。快其意所欲得。此人虽有百罪。弗法。即有避。因其事夷之。亦灭宗。以其故齐赵之郊盗贼不敢近广平。广平声为道不拾遗。上闻。迁为河内太守。\\
\hspace*{2em}素居广平时。皆知河内豪奸之家。及往。九月而至。令郡具私马五十匹。为驿自河内至长安。部吏如居广平时方略。捕郡中豪猾。郡中豪猾相连坐千馀家。上书请。大者至族。小者乃死。家尽没入偿臧。奏行不过二三日。得可事。论报。至流血十馀里。河内皆怪其奏。以为神速。尽十二月。郡中毋声。毋敢夜行。野无犬吠之盗。其颇不得。失之旁郡国。黎来。会春。温舒顿足叹曰。﹁嗟乎。令冬月益展一月。足吾事矣。﹂其好杀伐行威不爱人如此。天子闻之。以为能。迁为中尉。其治复放河内。徙诸名祸猾吏与从事。河内则杨皆。麻戊。关中杨赣。成信等。义纵为内史。惮未敢恣治。及纵死。张汤败後。徙为廷尉。而尹齐为中尉。\\
\hspace*{2em}尹齐者。东郡茌平人。以刀笔稍迁至御史。事张汤。张汤数称以为廉武。使督盗贼。所斩伐不避贵戚。迁为关内都尉。声甚於宁成。上以为能。迁为中尉。吏民益凋敝。尹齐木彊少文。豪恶吏伏匿而善吏不能为治。以故事多废。抵罪。上复徙温舒为中尉。而杨仆以严酷为主爵都尉。\\
\hspace*{2em}杨仆者。宜阳人也。以千夫为吏。河南守案举以为能。迁为御史。使督盗贼关东。治放尹齐。以为敢挚行。稍迁至主爵都尉。列九卿。天子以为能。南越反。拜为楼船将军。有功。封将梁侯。为荀彘所缚。居久之。病死。\\
\hspace*{2em}而温舒复为中尉。为人少文。居廷惛惛不辩。至於中尉则心开。督盗贼。素习关中俗。知豪恶吏。豪恶吏尽复为用。为方略。吏苛察。盗贼恶少年投缿购告言奸。置伯格长以牧司奸盗贼。温舒为人谩。善事有埶者。即无埶者。视之如奴。有埶家。虽有奸如山。弗犯。无埶者。贵戚必侵辱。舞文巧诋下户之猾。以焄大豪。其治中尉如此。奸猾穷治。大抵尽靡烂狱中。行论无出者。其爪牙吏虎而冠。於是中尉部中中猾以下皆伏。有势者为游声誉。称治。治数岁。其吏多以权富。\\
\hspace*{2em}温舒击东越还。议有不中意者。坐小法抵罪免。是时天子方欲作通天台而未有人。温舒请覆中尉脱卒。得数万人作。上说。拜为少府。徙为右内史。治如其故。奸邪少禁。坐法失官。复为右辅。行中尉事。如故操。\\
\hspace*{2em}岁馀。会宛军发。诏徵豪吏。温舒匿其吏华成。及人有变告温舒受员骑钱。他奸利事。罪至族。自杀。其时两弟及两婚家亦各自坐他罪而族。光禄徐自为曰。﹁悲夫。夫古有三族。而王温舒罪至同时而五族乎。﹂\\
\hspace*{2em}温舒死。家直累千金。後数岁。尹齐亦以淮阳都尉病死。家直不满五十金。所诛灭淮阳甚多。及死。仇家欲烧其尸。尸亡去归葬。\\
\hspace*{2em}自温舒等以恶为治。而郡守。都尉。诸侯二千石欲为治者。其治大抵尽放温舒。而吏民益轻犯法。盗贼滋起。南阳有梅免。白政。楚有殷中。杜少。齐有徐勃。燕赵之间有坚卢。范生之属。大群至数千人。擅自号。攻城邑。取库兵。释死罪。缚辱郡太守。都尉。杀二千石。为檄告县趣具食。小群以百数。掠卤乡里者。不可胜数也。於是天子始使御史中丞。丞相长史督之。犹弗能禁也。乃使光禄大夫范昆。诸辅都尉及故九卿张德等衣绣衣。持节。虎符发兵以兴击。斩首大部或至万馀级。及以法诛通饮食。坐连诸郡。甚者数千人。数岁。乃颇得其渠率。散卒失亡。复聚党阻山川者。往往而群居。无可柰何。於是作﹁沈命法﹂。曰群盗起不发觉。发觉而捕弗满品者。二千石以下至小吏主者皆死。其後小吏畏诛。虽有盗不敢发。恐不能得。坐课累府。府亦使其不言。故盗贼浸多。上下相为匿。以文辞避法焉。\\
\hspace*{2em}减宣者。杨人也。以佐史无害给事河东守府。卫将军青使买马河东。见宣无害。言上。徵为大厩丞。官事辨。稍迁至御史及中丞。使治主父偃及治淮南反狱。所以微文深诋。杀者甚众。称为敢决疑。数废数起。为御史及中丞者几二十岁。王温舒免中尉。而宣为左内史。其治米盐。事大小皆关其手。自部署县名曹实物。官吏令丞不得擅摇。痛以重法绳之。居官数年。一切郡中为小治辨。然独宣以小致大。能因力行之。难以为经。中废。为右扶风。坐怨成信。信亡藏上林中。宣使郿令格杀信。吏卒格信时。射中上林苑门。宣下吏诋罪。以为大逆。当族。自杀。而杜周任用。\\
\hspace*{2em}杜周者。南阳杜衍人。义纵为南阳守。以为爪牙。举为廷尉史。事张汤。汤数言其无害。至御史。使案边失亡。所论杀甚众。奏事中上意。任用。与减宣相编。更为中丞十馀岁。\\
\hspace*{2em}其治与宣相放。然重迟。外宽。内深次骨。宣为左内史。周为廷尉。其治大放张汤而善候伺。上所欲挤者。因而陷之。上所欲释者。久系待问而微见其冤状。客有让周曰。﹁君为天子决平。不循三尺法。专以人主意指为狱。狱者固如是乎。﹂周曰。﹁三尺安出哉。前主所是著为律。後主所是疏为令。当时为是。何古之法乎。﹂\\
\hspace*{2em}至周为廷尉。诏狱亦益多矣。二千石系者新故相因。不减百馀人。郡吏大府举之廷尉。一岁至千馀章。章大者连逮证案数百。小者数十人。远者数千。近者数百里。会狱。吏因责如章告劾。不服。以笞掠定之。於是闻有逮皆亡匿。狱久者至更数赦十有馀岁而相告言。大抵尽诋以不道以上。廷尉及中都官诏狱逮至六七万人。吏所增加十万馀人。\\
\hspace*{2em}周中废。後为执金吾。逐盗。捕治桑弘羊。卫皇后昆弟子刻深。天子以为尽力无私。迁为御史大夫。家两子。夹河为守。其治暴酷皆甚於王温舒等矣。杜周初徵为廷史。有一马。且不全。及身久任事。至三公列。子孙尊官。家訾累数巨万矣。\\
\hspace*{2em}太史公曰。自郅都。杜周十人者。此皆以酷烈为声。然郅都伉直。引是非。争天下大体。张汤以知阴阳。人主与俱上下。时数辩当否。国家赖其便。赵禹时据法守正。杜周从谀。以少言为重。自张汤死後。网密。多诋严。官事浸以秏废。九卿碌碌奉其官。救过不赡。何暇论绳墨之外乎。然此十人中。其廉者足以为仪表。其污者足以为戒。方略教导。禁奸止邪。一切亦皆彬彬质有其文武焉。虽惨酷。斯称其位矣。至若蜀守冯当暴挫。广汉李贞擅磔人。东郡弥仆锯项。天水骆璧推咸。河东褚广妄杀。京兆无忌。冯翊殷周蝮鸷。水衡阎奉朴击卖请。何足数哉。何足数哉。\\

% ===== 索隐述赞（司马贞《史记索隐》）· 统一样式 =====
% 题头「索隐述赞曰」五字并入反白框；正文每四句一列，定宽盒 \hfil 居列中；
% \fontsize{24}{41} 列距对齐全书界行（121+41k）；\clearpage 使述赞独占一页。
\clearpage
\shuzanhead{索隐述赞曰}\linebreak
{\shuzanbody\fontsize{24}{41}\selectfont
\hbox to \hsize{\hfil 太上失德。法令滋起。破觚为圆。禁暴不止。\hfil}\linebreak
\hbox to \hsize{\hfil 奸伪斯炽。惨酷爰始。乳兽扬威。苍鹰侧视。\hfil}\linebreak
\hbox to \hsize{\hfil 舞文巧诋。怀生何恃。\hfil}%
\par}
\end{document}
